\chapter*{ملخص البحث}
\addcontentsline{toc}{chapter}{ملخص البحث}

\begin{Arabic}
\large
في إطار تحديث ورقمنة الإجراءات الإدارية لفائدة الفاعلين الاقتصاديين في مجال التجارة الخارجية بالمغرب، تتولى شركة \textbf{PortNet S.A.} (الشباك الموحد الوطني) إدارة طلبات الاشتراك والولوج إلى المنظومات المعلوماتية لـ PortNet وإدارة الجمارك والضرائب المباشرة وغير المباشرة (نظام BADR). يشكل التدقيق اليدوي التقليدي لوثائق الاشتراك (النموذج J، التعريف الضريبي IF، الرسم المهني TP، شهادة التعريف الموحد ICE، شهادة التوطين، كشف الحساب البنكي RIB، بطاقة التعريف الوطنية CIN، والنظام الأساسي للشركة) عائقاً رئيسياً يتسبب في تأخير معالجة الطلبات، وارتفاع مخاطر التزوير، والأخطاء البشرية.

يقدم هذا التقرير الخاص بـ \textbf{مشروع نهاية السنة (PFA)} تصميم وبناء منصة إلكترونية ذكية ومؤتمتة للتحقق من صحة ومطابقة وثائق الاشتراك. تعتمد المنصة على بنية برمجية حديثة تفصل بين واجهة مستخدم تفاعلية (HTML5/CSS3/Vanilla JS) وخلفية نظام عالية الأداء باستعمال لغة Python وإطار العمل FastAPI وسيرفر Uvicorn.

تتجلى المساهمة التقنية الرئيسية لهذا العمل في ابتكار \textbf{محرك تصنيف والتحقق ذو طبقتين (Two-Layer Slot Classification Engine)}. يقضي هذا المحرك بشكل نهائي على مشكلة القراءات الخاطئة والتداخل بين الوثائق (*Cross-Slot Confusion*) عبر الدمج بين علامات الهيكلة الإيجابية الصارمة (رموز النماذج الرسمية \texttt{AAC140B} و \texttt{AAC060B} وشريط المقروئية الآلية \texttt{<<}) وقواعد الاستبعاد السلبية الإجبارية (*Negation Blockers*). علاوة على ذلك، تم تطوير خوارزميات استخراج متخصصة لمعالجة تعقيدات طبقات النصوص في ملفات PDF، مع تفعيل تدقيق تاريخ الصلاحية (قاعدة 3 أشهر للنموذج J)، والتحقق الرياضياتي باستخدام خوارزمية Modulo 97 (لأرقام ICE و RIB)، بالإضافة إلى مصفوفة المطابقة التقاطعية بين الوثائق مع دعم تصحيح الأخطاء المطبعية في أرقام بطاقات التعريف.

تم تقييم المنصة بشكل دقيق عبر مصفوفة اختبارات شاملة (8x8) شملت وثائق رقمية حقيقية للشركات عبر جميع منافذ الرفع. وأظهرت النتائج التجريبية تحقيق نسبة نجاح \textbf{100% (64/64 اختباراً ناجحاً)} بدون أي إنذار خاطئ، وبزمن استجابة متوسط يقل عن 350 ميلي ثانية لكل وثيقة، مع حفظ البيانات في قاعدة بيانات PostgreSQL بطريقة منظمة وآمنة.

\vspace{0.5cm}
\noindent\textbf{الكلمات المفتاحية:} بورنيت (PortNet S.A.)، الجمارك (نظام BADR)، التدقيق الآلي للوثائق، استخراج البيانات، FastAPI، محرك التصنيف ذو الطبقتين، خوارزمية Modulo 97، مصفوفة التناقضات المتقاطعة، PostgreSQL.
\end{Arabic}
